\section{Representation facts and normalization}\label{sec:prelim}
\label{ac:part:1}

A permutation sends each input position to its output, and a product
\(gh\) applies \(h\) first. Thus convolution is the law of composition,
and
\[
 \widehat{\mu*\nu}(\lambda)
 =\widehat\mu(\lambda)\widehat\nu(\lambda),
 \qquad
 \widehat\mu(\lambda)=\sum_g\mu(g)\rho_\lambda(g).
\]
We use ordinary, unnormalized matrix traces:
\(\|A\|_p^p=\Tr|A|^p\), where \(|A|=(A^*A)^{1/2}\).

The group acts on slots, meaning positions for cards. A random bijection
between two slot sets of the same size becomes a random permutation
after identifying each set with a fixed reference set. Changing either
identification multiplies its average in a unitary representation by
unitaries on the two sides, so its singular values do not change. We
will use this convention when prescribed paths leave different free
slot sets at successive layers.

All representations are over \(\mathbb C\). For a partition
\(\lambda\vdash l\), write \([\lambda]\) for the irreducible
representation of \(S_l\), \(D_\lambda\) for its dimension, and
\(F(\lambda)=\log D_\lambda\). The empty diagram, for \(l=0\), has
dimension one. All logarithms in the proof are natural. We use the
following classical representation facts.
\begin{itemize}
\item The dimension \(D_\lambda\) is the number of standard tableaux
of shape \(\lambda\), given by the hook-length formula. Conjugating
the diagram, denoted by \(\lambda'\), twists the representation by
sign. For \(0\le j\le l\), on restriction to \(S_{l-j}\) the multiplicity of \([\sigma]\)
is the number of standard tableaux of the skew shape
\(\lambda/\sigma\) when \(\sigma\subseteq\lambda\), and is zero
otherwise. These are the branching and tableau rules; see
Sagan and Stanley~\cite{Sagan2001,Stanley1999} and
Vershik--Okounkov~\cite[Theorem~5.8 and Corollary~7.1]{VershikOkounkov2005}.
In particular, the representation on placements of \(j\)
distinguishable cards in \(l\) slots contains \([\lambda]\) exactly
when \(\lambda_1\ge l-j\).

\item Pieri's rule says that inducing with a trivial representation
adds a horizontal strip, while inducing with a sign representation
adds a vertical strip~\cite{Sagan2001,Stanley1999}.

\item In ordinary Schur--Weyl duality, for integers \(p\ge1\) and \(t\ge0\), the \(S_t\)-types in
\((\mathbb C^p)^{\otimes t}\) are the partitions \(\sigma\vdash t\)
with at most \(p\) rows. Their multiplicity spaces are the polynomial
irreducibles of highest weight \(\sigma\) for the commuting linear
group action~\cite{Sagan2001,Stanley1999}. Section~\ref{sec:density}
uses this decomposition separately on two parity classes and records
the quantitative multiplicity and weight estimates there.
\end{itemize}

The induction separates representations close to a row or column from
those with exponentially large dimension. For large \(l\), put
\[
 k=\min(l-\lambda_1,l-\lambda'_1).
\]
If \(0<k\le l/3\), then
\begin{equation}
 D_\lambda\ge e^{-1}\binom{l}{k}.
 \label{ac:eq:1}
\end{equation}
Indeed, after transposing if necessary, the first row has length \(l-k\).
Let \(\mu\) be the diagram of \(k\) boxes below it. The hook formula gives
\[
 D_\lambda
 =\binom{l}{k}D_\mu
   \prod_{j=1}^{l-k}
   \left(1+\frac{\mu'_j}{l-k-j+1}\right)^{-1}.
\]
Only \(j\le k\) can contribute a nontrivial factor. Since
\(\sum_j\mu'_j=k\), the product is at least
\(\exp(-k/(l-2k))\ge e^{-1}\), proving \eqref{ac:eq:1}.

If \(k>l/3\), then
\[
 F(\lambda)\ge c_{\rm dim}l
\]
for an absolute \(c_{\rm dim}>0\) and all large \(l\). If both the first
row and first column have length below \(l/8\), every hook is shorter
than \(l/4\), and the hook formula with Stirling's bound gives this
conclusion. Otherwise transpose if needed so that the first row has
length \(v\in[l/8,2l/3]\). Retain that row and a subdiagram of
\(\lfloor v/4\rfloor\) boxes below it. The preceding large-row bound
makes the dimension of this retained shape exponential in \(v\), and
hence in \(l\). Branching shows that \(D_\lambda\) is at least that
dimension.

Reducing \(c_{\rm dim}\) if necessary, these estimates imply
\(F(\lambda)\ge c_{\rm dim}k\) for every \(k\). Conversely the placement
criterion, with a sign twist if necessary, places \([\lambda]\) inside
a representation of dimension \((l)_k\le l^k\). Thus
\[
 c_{\rm dim}k\le F(\lambda)\le k\log l
\]
for large \(l\). These are the bounds used to enter the sparse range;
\eqref{ac:eq:1} also controls the sum over representations in the final
mixing argument.
